\documentclass[12pt]{article}
\usepackage[utf8]{inputenc}
\usepackage[T1]{fontenc}
\usepackage[margin=1in]{geometry}
\usepackage{setspace}
\usepackage{xcolor}
\usepackage[colorlinks=true, urlcolor=blue!60!black, linkcolor=blue!60!black]{hyperref}
\setlength{\parindent}{2em}
\setlength{\parskip}{1.25em}
\renewcommand{\baselinestretch}{1.0}
\begin{document}
\begin{center}
{\LARGE \textit{United States v. Morrison}, 529 U.S. 598 (2000)\par}
\vspace{0.8em}
{\large Case Brief\par}
\end{center}
\vspace{1em}

\section{Memory Jogger}
Applying \textit{Lopez} and the Fourteenth Amendment's state-action requirement, the Supreme Court invalidated the federal civil remedy of the Violence Against Women Act of 1994, holding that Congress may not regulate noneconomic violent crime based solely on aggregate economic effects---even with voluminous congressional findings---nor remedy state bias by creating a nationwide private cause of action against private individuals (\href{https://www.courtlistener.com/opinion/118363/united-states-v-morrison/}{CourtListener}; \href{https://www.law.cornell.edu/supct/html/99-5.ZO.html}{Cornell LII}).

\section{Detailed Case Facts}
In the fall of 1994, Christy Brzonkala enrolled at Virginia Polytechnic Institute (Virginia Tech) (\href{https://www.law.cornell.edu/supct/html/99-5.ZO.html}{Majority Opinion}). In September 1994, she met Antonio Morrison and James Crawford, fellow students and members of Virginia Tech's varsity football team. Brzonkala alleged that within 30 minutes of meeting them, Morrison and Crawford assaulted and repeatedly raped her. Afterward, Morrison allegedly told her, ``You better not have any \ldots{} diseases,'' and in the following months allegedly boasted in the dormitory dining room that he ``like[d] to get girls drunk and \ldots{}'' in crude, vulgar terms (\href{https://www.law.cornell.edu/supct/html/99-5.ZO.html}{Majority Opinion}). Brzonkala suffered severe emotional distress and depression, sought treatment and antidepressant medication from a university psychiatrist, stopped attending classes, and withdrew from Virginia Tech.

In early 1995, Brzonkala filed a complaint against Morrison and Crawford under Virginia Tech's Sexual Assault Policy. At the university hearing, Morrison admitted having sexual contact with Brzonkala despite her twice telling him ``no'' (\href{https://www.law.cornell.edu/supct/html/99-5.ZO.html}{Majority Opinion}). Virginia Tech's Judicial Committee found insufficient evidence to punish Crawford, found Morrison guilty of sexual assault, and imposed an immediate two-semester suspension, which the dean of students upheld. In July 1995, however, university officials informed Brzonkala that a second hearing was required because the Sexual Assault Policy had not been widely circulated to students. After a second hearing under Virginia Tech's Abusive Conduct Policy, the Judicial Committee again found Morrison guilty and imposed a two-semester suspension, but without explanation changed the description of his offense from ``sexual assault'' to ``using abusive language.'' On August 21, 1995, Virginia Tech's senior vice president and provost set aside Morrison's suspension as ``excessive'' compared with other Abusive Conduct Policy cases, without notifying Brzonkala. After reading in a newspaper that Morrison would return for the fall 1995 semester, Brzonkala dropped out of Virginia Tech (\href{https://www.law.cornell.edu/supct/html/99-5.ZO.html}{Majority Opinion}; \href{https://www.oyez.org/cases/1999/99-5}{Oyez}).

At issue was 42 U.S.C. \S 13981, enacted as \S 40302 of the Violence Against Women Act of 1994 (VAWA), 108 Stat. 1941--1942. Section 13981 declared a right to be free from ``crimes of violence motivated by gender'' and created a federal civil cause of action for compensatory and punitive damages, injunctive relief, and declaratory relief against perpetrators, regardless of whether criminal charges were filed (\href{https://www.law.cornell.edu/supct/html/99-5.ZO.html}{Majority Opinion}). Congress explicitly invoked both Article I, \S 8 (the Commerce Clause) and Section 5 of the Fourteenth Amendment as constitutional authority for \S 13981 (\S 13981(a)), supported by four years of congressional hearings, eight committee reports, and task force reports from 21 States documenting both the economic toll of violence against women and gender bias in state justice systems (\href{https://www.law.cornell.edu/supct/html/99-5.ZD.html}{Souter Dissent}).

\section{Procedural History}
In December 1995, Brzonkala sued Morrison, Crawford, and Virginia Tech in the United States District Court for the Western District of Virginia, alleging that Morrison and Crawford violated 42 U.S.C. \S 13981 and that Virginia Tech's handling of her complaint violated Title IX of the Education Amendments of 1972, 20 U.S.C. \S\S 1681--1688 (\href{https://www.law.cornell.edu/supct/html/99-5.ZO.html}{Majority Opinion}). Morrison and Crawford moved to dismiss for failure to state a claim and on the ground that \S 13981 was unconstitutional. The United States intervened to defend the statute's constitutionality.

The District Court dismissed Brzonkala's Title IX claims against Virginia Tech, 935 F. Supp. 772 (W.D. Va. 1996), and held that while her complaint stated a claim against Morrison and Crawford under \S 13981, Congress lacked constitutional authority to enact \S 13981 under either the Commerce Clause or Section 5 of the Fourteenth Amendment, 935 F. Supp. 779 (W.D. Va. 1996) (\href{https://www.law.cornell.edu/supct/html/99-5.ZO.html}{Majority Opinion}). A divided panel of the Fourth Circuit reversed, reinstating Brzonkala's \S 13981 claim and Title IX hostile environment claim, 132 F.3d 949 (4th Cir. 1997). The full Fourth Circuit vacated the panel opinion and reheard the case en banc; by a divided vote, the en banc court affirmed the District Court's holding that \S 13981 was unconstitutional, 169 F.3d 820 (4th Cir. 1999) (while remanding the Title IX hostile environment claim to be held in abeyance pending \textit{Davis v. Monroe County Bd. of Ed.}, 526 U.S. 629 (1999)). The Supreme Court granted certiorari in the consolidated petitions of the United States (No. 99-5) and Brzonkala (No. 99-29), 527 U.S. 1068 (1999), heard oral argument on January 11, 2000, and affirmed on May 15, 2000 (\href{https://www.oyez.org/cases/1999/99-5}{Oyez}; \href{https://www.law.cornell.edu/supct/html/99-5.ZS.html}{Syllabus}).

\section{Judicial Votes}
\textbf{Vote:} 5--4 affirming the en banc Fourth Circuit (\href{https://www.law.cornell.edu/supct/html/99-5.ZS.html}{Syllabus}; \href{https://www.courtlistener.com/opinion/118363/united-states-v-morrison/}{CourtListener}).
\begin{itemize}
\item \textbf{Majority Opinion:} Chief Justice Rehnquist, joined by Justices O'Connor, Scalia, Kennedy, and Thomas (\href{https://www.law.cornell.edu/supct/html/99-5.ZO.html}{Majority Opinion}).
\item \textbf{Concurring Opinion:} Justice Thomas (\href{https://www.law.cornell.edu/supct/html/99-5.ZC.html}{Thomas Concurrence}).
\item \textbf{Dissenting Opinions:} Justice Souter, joined by Justices Stevens, Ginsburg, and Breyer (\href{https://www.law.cornell.edu/supct/html/99-5.ZD.html}{Souter Dissent}); Justice Breyer, joined by Justice Stevens, and joined by Justices Souter and Ginsburg as to Part I-A (\href{https://www.law.cornell.edu/supct/html/99-5.ZD1.html}{Breyer Dissent}).
\end{itemize}

\section{Holding}
\textbf{Judgment:} Affirmed. Section 40302 of the Violence Against Women Act of 1994, 42 U.S.C. \S 13981, cannot be sustained under either the Commerce Clause (U.S. Const., Art. I, \S 8, cl. 3) or Section 5 of the Fourteenth Amendment (\href{https://www.law.cornell.edu/supct/html/99-5.ZO.html}{Majority Opinion}).

\textbf{Rule Announced:} (1) Under the Commerce Clause, Congress may not regulate noneconomic, violent criminal conduct based solely on that conduct's aggregate effect on interstate commerce, even where Congress compiles extensive legislative findings relying on a but-for causal chain to national economic impacts. (2) Under Section 5 of the Fourteenth Amendment, Congress's remedial enforcement power is limited to state action and must exhibit ``congruence and proportionality'' to a constitutional violation by the States; a nationwide civil remedy directed solely against private individuals rather than state officials exceeds Section 5 (\href{https://www.law.cornell.edu/supct/html/99-5.ZO.html}{Majority Opinion}).

\textbf{Holding vs. Dicta:} The Court explicitly held that \S 13981 is unconstitutional under both the Commerce Clause and Section 5, while leaving intact other VAWA provisions such as 18 U.S.C. \S 2261(a)(1) (which criminalizes interstate travel to commit domestic violence and was cited approvingly in footnote 5 as a valid category-one channels regulation) and VAWA's funding programs (\href{https://www.law.cornell.edu/supct/html/99-5.ZO.html}{Majority Opinion}). Notably, the majority expressly declined to ``adopt a categorical rule against aggregating the effects of any noneconomic activity''---a carefully reserved boundary that keeps broader categorical aggregation questions in dicta rather than strict holding.

\section{Analysis of Opinions}

\subsection{Majority Opinion (Chief Justice Rehnquist)}
Chief Justice Rehnquist analyzed \S 13981 under both asserted constitutional grants of power, beginning with the Commerce Clause and applying the four-factor framework distilled from \textit{United States v. Lopez}, 514 U.S. 549 (1995) (\href{https://www.law.cornell.edu/supct/html/99-5.ZO.html}{Majority Opinion}). Because petitioners relied solely on \textit{Lopez}'s third category (activities substantially affecting interstate commerce), Rehnquist evaluated: (1) the economic or noneconomic nature of the regulated activity; (2) the presence of an express jurisdictional element; (3) the existence of congressional findings; and (4) the degree of attenuation between the regulated conduct and interstate commerce.

First, Rehnquist emphasized that gender-motivated crimes of violence ``are not, in any sense of the phrase, economic activity,'' noting that the Court had historically upheld Commerce Clause aggregation only where the regulated intrastate activity was economic in nature (\href{https://www.law.cornell.edu/supct/html/99-5.ZO.html}{Majority Opinion}). Second, like the statute in \textit{Lopez}, \S 13981 contained no express jurisdictional element limiting its reach to conduct connected to interstate commerce. Third, although Congress had amassed voluminous findings (unlike in \textit{Lopez}), Rehnquist declared that whether operations substantially affect interstate commerce is ultimately a judicial question, and Congress's findings were ``substantially weakened'' because they relied on the exact ``but-for causal chain'' from violent crime to employment, travel, medical costs, and national productivity that \textit{Lopez} had rejected. Accepting that logic would empower Congress to regulate murder, all violent crime, and even marriage, divorce, and childrearing---erasing the line between ``what is truly national and what is truly local'' and usurping the States' core police power.

Turning to Section 5 of the Fourteenth Amendment, Rehnquist invoked \textit{United States v. Harris}, 106 U.S. 629 (1883), and the \textit{Civil Rights Cases}, 109 U.S. 3 (1883), to reaffirm the ``time-honored principle'' that the Fourteenth Amendment prohibits only state action, not private conduct (\href{https://www.law.cornell.edu/supct/html/99-5.ZO.html}{Majority Opinion}). Rehnquist dismissed petitioners' reliance on \textit{United States v. Guest}, 383 U.S. 745 (1966), and \textit{District of Columbia v. Carter}, 409 U.S. 418 (1973), as mere dicta. Even assuming pervasive gender bias in state justice systems, \S 13981 failed the ``congruence and proportionality'' test of \textit{City of Boerne v. Flores}, 521 U.S. 507 (1997): it was directed not at any State or state official (visiting no consequence on any Virginia official), unlike remedies upheld in \textit{Katzenbach v. Morgan}, 384 U.S. 641 (1966), and \textit{South Carolina v. Katzenbach}, 383 U.S. 301 (1966), and it applied uniformly nationwide even though Congress's findings did not show bias in all or even most States.

\subsection{Concurring Opinion (Justice Thomas)}
In a concise one-paragraph concurrence, Justice Thomas joined the majority opinion in full as a correct application of \textit{Lopez}, but reiterated his originalist objection from \textit{Lopez} that the ``substantial effects'' test itself is ``rootless and malleable'' and inconsistent with the original understanding of the Commerce Clause (\href{https://www.law.cornell.edu/supct/html/99-5.ZC.html}{Thomas Concurrence}). Thomas warned that until the Court replaces the substantial effects test with an originalist standard, Congress will continue appropriating state police powers under the guise of regulating commerce.

\subsection{Dissenting Opinion (Justice Souter, joined by Justices Stevens, Ginsburg, and Breyer)}
Justice Souter argued that the majority's holding was irreconcilable with settled Commerce Clause precedent from \textit{Wickard v. Filburn} through \textit{Heart of Atlanta Motel, Inc. v. United States}, 379 U.S. 241 (1964), and \textit{Katzenbach v. McClung}, 379 U.S. 294 (1964) (\href{https://www.law.cornell.edu/supct/html/99-5.ZD.html}{Souter Dissent}). Unlike \textit{Lopez}, Congress had assembled a ``mountain of data'' over four years of hearings---including estimates that domestic violence costs the nation \$3 billion to \$10 billion annually and that nearly 50 percent of rape victims lose or quit their jobs---amply satisfying rational basis review.

Souter accused the majority of replacing factual rational basis scrutiny with a new, judicially contrived categorical exclusion for ``noneconomic'' and ``traditional state'' subjects, reviving the pre-1937 formalism of \textit{Carter Coal} and \textit{Schechter Poultry}---this time in service of a judicial vision of federalism rather than laissez-faire economics (\href{https://www.law.cornell.edu/supct/html/99-5.ZD.html}{Souter Dissent}). Citing \textit{Gibbons v. Ogden} and \textit{Garcia v. San Antonio Metropolitan Transit Authority}, 469 U.S. 528 (1985), Souter maintained that the Constitution entrusts the federal-state balance within the commerce sphere to the national political process, noting the irony that 36 States and Puerto Rico filed an amicus brief supporting \S 13981 while only one State opposed it: Antonio Morrison had ``won the states' rights plea against the states themselves.'' Because Souter would have upheld \S 13981 under the Commerce Clause, he did not reach Section 5.

\subsection{Dissenting Opinion (Justice Breyer, joined by Justice Stevens, and joined by Justices Souter and Ginsburg as to Part I-A)}
In Part I-A (joined by all four dissenters), Justice Breyer critiqued the majority's ``economic/noneconomic'' boundary as unworkable and arbitrary (\href{https://www.law.cornell.edu/supct/html/99-5.ZD1.html}{Breyer Dissent}). He asked whether a street-corner mugger engages in economic activity, noted that the Court allows aggregation of noneconomic discrimination at commercial establishments (\textit{Heart of Atlanta}) and noneconomic conduct inside broader economic schemes, and observed that Congress can frequently bypass subject-matter limits simply by adding a jurisdictional hook tied to goods that crossed state lines. Because modern economic integration means almost any local activity can affect interstate commerce in the aggregate, Congress---not the judiciary---must bear primary responsibility for striking the federal-state balance.

In Part I-B (joined by Justice Stevens), Breyer suggested that if judicial limits are to be applied at all, a procedural ``hard look'' approach focusing on the thoroughness of Congress's legislative deliberation and state-federal cooperation (here, support from 38 state Attorneys General and an explicit carve-out for family law in \S 13981(e)(4)) would be far more workable than rigid substantive categories (\href{https://www.law.cornell.edu/supct/html/99-5.ZD1.html}{Breyer Dissent}). In Part II (also joined by Justice Stevens), Breyer expressed serious doubt about the majority's Section 5 analysis: unlike the statutes in \textit{Harris} and the \textit{Civil Rights Cases}, VAWA responded directly to documented state-actor bias in state courts, and providing a federal civil substitute against private wrongdoers for conduct already unlawful under state law was a minimally intrusive, congruent, and proportionate remedy.

\section{Examples: Future Applications}
\begin{enumerate}
\item \textbf{Same-Side (Invalid under \textit{Morrison}): Federal Civil Tort Remedy for Local Hate Crimes Without a Jurisdictional Element.} Suppose Congress creates a federal civil damages cause of action for any victim of a racially or religiously motivated assault, relying solely on congressional findings that hate crimes deter interstate travel and reduce workforce productivity, with no jurisdictional element and no link to state actors. Under \textit{Morrison}, the regulated conduct is noneconomic violent crime, aggregate economic findings cannot overcome that defect under the Commerce Clause, and a purely private civil remedy nationwide cannot be sustained under Section 5 of the Fourteenth Amendment.
\item \textbf{Same-Side (Invalid under \textit{Morrison}): Federal Anti-Arson Statute Applied to an Owner-Occupied Private Residence.} If Congress prosecutes an individual under a federal arson statute for burning down a purely private, owner-occupied home that is not rented or used in any commercial activity, \textit{Morrison} and \textit{Lopez} bar stretching the Commerce Clause via aggregate insurance or utility-connection reasoning (consistent with the Court's statutory avoidance ruling in \textit{Jones v. United States}, 529 U.S. 848 (2000), decided one week after \textit{Morrison}).
\item \textbf{Opposite-Side (Valid under \textit{Morrison}): Federal Criminalization of Interstate Domestic Violence (18 U.S.C. \S 2261(a)(1)).} A federal statute punishing any person who crosses a state line with the intent to injure or harass an intimate partner and commits a crime of violence in the course of that travel is valid under \textit{Lopez}'s first category (regulating the channels of interstate commerce), as Chief Justice Rehnquist expressly noted in footnote 5 of \textit{Morrison} (\href{https://www.law.cornell.edu/supct/html/99-5.ZO.html}{Majority Opinion}).
\item \textbf{Opposite-Side (Valid under \textit{Morrison}): Section 5 Legislation Targeting State Officials' Discriminatory Policing.} Suppose Congress, backed by a documented record of systematic equal-protection violations by specific state police departments, enacts a statute authorizing federal injunctive suits and damages directly against those state agencies and officers for refusing to investigate sexual assaults. Because the statute targets \textit{state action} directly rather than private perpetrators and is tailored to documented constitutional violations, it satisfies both the state-action requirement and \textit{Boerne}'s congruence-and-proportionality standard.
\item \textbf{Fence-Sitter (Close Call): Federal Protection of Purely Intrastate Endangered Species on Private Land.} Suppose Congress prohibits a private landowner from clearing brush on residential property that serves as the sole habitat for a tiny, noncommercial endangered species living entirely within one State. Under \textit{Morrison}, the landowner's noncommercial conduct looks local and noneconomic, suggesting invalidity. Yet under \textit{Lopez}'s ``essential part of a larger regulation of economic activity'' exception---and later \href{https://www.law.cornell.edu/supremecourt/text/03-1454}{\textit{Gonzales v. Raich}} (2005)---defenders argue that protecting biodiversity is integral to the Endangered Species Act's comprehensive regulatory scheme governing interstate markets in wildlife, medicine, and ecosystems, making it a perennial constitutional boundary dispute.
\end{enumerate}

\section{Critique}
\textit{United States v. Morrison} proved even more consequential---and controversial---than \textit{Lopez} because it demonstrated that extensive congressional fact-finding could not rescue a statute regulating noneconomic conduct, while simultaneously tightening the state-action and congruence-and-proportionality constraints of Section 5 (\href{https://www.law.cornell.edu/constitution-conan/amendment-10/the-commerce-clause-and-the-tenth-amendment}{Cornell LII Constitution Annotated}).

From an \textbf{originalist and textualist perspective}, scholars such as Randy Barnett praised \textit{Morrison} for confirming that \textit{Lopez} was not an isolated fluke and for defending the Constitution's structural enumeration of powers against a limitless ``but-for'' aggregation theory (\href{https://scholarship.law.georgetown.edu/cgi/viewcontent.cgi?article=1856&context=facpub}{Barnett, 72 U. Colo. L. Rev. 1275 (2002)}). Yet originalists remained divided: like Justice Thomas, strict originalists criticized the majority for leaving the post-1937 ``substantial effects'' framework intact, while some historical scholars questioned the majority's uncritical reliance on the 1883 \textit{Civil Rights Cases} and \textit{Harris} as faithful reflections of the Fourteenth Amendment's original Reconstruction meaning.

From a \textbf{progressive, feminist, and doctrinal perspective}, scholars leveled sharp critiques against both halves of the opinion. Robert Post and Reva Siegel argued that \textit{Morrison} combined a constricted Commerce Clause with an aggressive judicial monopoly over Fourteenth Amendment enforcement, insulating gender-based violence and family-sphere inequality as ``local'' and ``noneconomic'' despite Congress's explicit constitutional mandate to enforce equal protection (\href{https://law.yale.edu/sites/default/files/documents/pdf/Faculty/Siegel_EqualProtectionByLaw.pdf}{Post \& Siegel, 110 Yale L.J. 441 (2000)}). Kermit Roosevelt III similarly argued that \textit{Morrison} was wrong on Section 5 because it misapplied the state-action doctrine to invalidate a remedial statute aimed directly at documented state-court failures to provide equal protection (\href{https://scholarship.law.cornell.edu/cgi/viewcontent.cgi?article=4658&context=clr}{Roosevelt, 100 Cornell L. Rev. (2015)}). Erwin Chemerinsky criticized the Court's federalism revival as resting on unexamined assumptions that judicial line-drawing protects liberty better than cooperative federal-state civil rights remedies (\href{http://www.stanfordlawreview.org/wp-content/uploads/sites/3/2010/04/chemerinsky.pdf}{Chemerinsky, 58 Stan. L. Rev. 1763 (2006)}).

At its core, \textit{Morrison} exposes the tension between \textbf{structural dual federalism}---insisting on judicially policed boundaries reserving traditional criminal and family law to the States---and \textbf{national civil rights enforcement and cooperative federalism}, where Congress and a supermajority of States seek federal remedies to overcome persistent, systemic failures in local justice systems.

\section{Key Quotations}
\begin{quote}
``We accordingly reject the argument that Congress may regulate noneconomic, violent criminal conduct based solely on that conduct's aggregate effect on interstate commerce. The Constitution requires a distinction between what is truly national and what is truly local.''
--- Chief Justice Rehnquist, Opinion of the Court (\href{https://www.law.cornell.edu/supct/html/99-5.ZO.html}{Cornell LII}; \href{https://www.courtlistener.com/opinion/118363/united-states-v-morrison/}{CourtListener}).
\end{quote}

\begin{quote}
``If the allegations here are true, no civilized system of justice could fail to provide her a remedy for the conduct of respondent Morrison. But under our federal system that remedy must be provided by the Commonwealth of Virginia, and not by the United States.''
--- Chief Justice Rehnquist, Opinion of the Court (\href{https://www.law.cornell.edu/supct/html/99-5.ZO.html}{Cornell LII}; \href{https://www.courtlistener.com/opinion/118363/united-states-v-morrison/}{CourtListener}).
\end{quote}

\begin{quote}
``Until this Court replaces its existing Commerce Clause jurisprudence with a standard more consistent with the original understanding, we will continue to see Congress appropriating state police powers under the guise of regulating commerce.''
--- Justice Thomas, Concurring (\href{https://www.law.cornell.edu/supct/html/99-5.ZC.html}{Cornell LII}).
\end{quote}

\begin{quote}
``One obvious difference from United States v. Lopez, 514 U.S. 549 (1995), is the mountain of data assembled by Congress, here showing the effects of violence against women on interstate commerce.''
--- Justice Souter, Dissenting, joined by Justices Stevens, Ginsburg, and Breyer (\href{https://www.law.cornell.edu/supct/html/99-5.ZD.html}{Cornell LII}).
\end{quote}

\begin{quote}
``Since judges cannot change the world, the `defect' means that, within the bounds of the rational, Congress, not the courts, must remain primarily responsible for striking the appropriate state/federal balance.''
--- Justice Breyer, Dissenting, joined by Justice Stevens, and joined by Justices Souter and Ginsburg as to Part I-A (\href{https://www.law.cornell.edu/supct/html/99-5.ZD1.html}{Cornell LII}).
\end{quote}

\end{document}
